\scanpage{24}
- \mspage{18} -
qu'il est divisible par $c_4^3-c_6^2$ dans $\mathbb Z/(2)[c_4,c_6,\Delta]$ (resp. $\mathbb Z/(3)[c_4,c_6,\Delta]$). Cela résulte de la structure de l'anneau des formes modulaires de caractéristique 2 (resp. 3).

\uline{Proposition} 7.2.\quad(I) \uline{L'anneau des formes modulaires de caractéristique} 2 \uline{est} $\mathbb Z/(2)[a_1,\Delta]$, \uline{et} $c_4=a_1^4$, $c_6=a_1^6$

(II) \uline{L'anneau des formes modulaires de caractéristique} 3 \uline{est} $\mathbb Z/(3)[b_2,\Delta]$, et $c_4=b_2^2$, $c_6=-b_2^3$.

Tout d'abord, il résulte de $(1.6)_1$ et $(1.7)_2$ que $a_1$ (resp. $b_2$) est une forme modulaire de caractéristique 2 (resp. 3).

En caractéristique 2, cherchons tout d'abord les formes modulaires relatives aux seules courbes telles que $a_1$ soit inversible. Ces courbes se mettent sous forme (cf. (4.1))
\[
y^2+a_1xy=x^3+a_2x^2+a_6\qquad\pi=\frac{dx}{a_1x}
\]
et, compte tenu de (4.3), on voit que ces formes s'écrivent
\[
a_1^{-n}\ P(a_1,a_6)\qquad\text{pour }P\text{ isobare, ou encore}
\]
\[
a_1^{-m}\ Q(a_1,\Delta)
\]
et cela garde un sens pour $a_1=0$, $\Delta\ne0$ si et seulement si on avait en fait affaire à un polynome en $a_1$ et $\Delta$.

En caractéristique 3, raisonnant de même d'abord sur les courbes telles que $c_4$ soit inversible, et remarquant que $d_6=0$ équivaut alors à $a_4=0$ dans (3.1), on voit que toute forme modu-
\margin{23}
